% =====================================================================
%  软共线有效理论（SCET）课前预习讲义
%  课程：SCET / Soft-Collinear Effective Theory
%  授课：Matthias Neubert
%  说明：本文件为自包含（self-contained）LaTeX 源文件，
%        请使用 XeLaTeX（推荐）或 LuaLaTeX 编译两次以生成正确目录。
%        编译命令示例： xelatex 03-scet-matthias-neubert.tex
% =====================================================================
\documentclass[UTF8,a4paper,11pt]{ctexart}

\usepackage{amsmath}
\usepackage{amssymb}
\usepackage{amsthm}
\usepackage{geometry}
\usepackage{booktabs}
\usepackage{enumitem}
\usepackage{mathtools}
\usepackage{bm}
\usepackage{hyperref}

\geometry{left=2.4cm,right=2.4cm,top=2.6cm,bottom=2.6cm}

\hypersetup{
  colorlinks=true,
  linkcolor=blue,
  citecolor=blue,
  urlcolor=magenta,
  pdftitle={软共线有效理论（SCET）课前预习讲义},
  pdfauthor={课前自学材料},
  pdfsubject={Soft-Collinear Effective Theory, Matthias Neubert}
}

% ---------------------------------------------------------------------
%  定理类环境（环境名用 ASCII，显示名用中文）
% ---------------------------------------------------------------------
\theoremstyle{definition}
\newtheorem{definition}{定义}[section]
\newtheorem{example}{例}[section]
\newtheorem{exercise}{练习}[section]

\theoremstyle{remark}
\newtheorem{remark}{注}[section]
\newtheorem{keypoint}{要点}[section]

\theoremstyle{plain}
\newtheorem{proposition}{命题}[section]

% ---------------------------------------------------------------------
%  常用记号宏（不依赖 slashed 宏包，用负核距自制斜线）
% ---------------------------------------------------------------------
\newcommand{\nsl}{n\mkern-10.5mu/\mkern1.5mu}
\newcommand{\nbsl}{\bar{n}\mkern-10.5mu/\mkern1.5mu}
\newcommand{\Dsl}{D\mkern-11.5mu/}
\newcommand{\psl}{p\mkern-9.5mu/}
\newcommand{\ksl}{k\mkern-9.0mu/}
\newcommand{\lsl}{\ell\mkern-8.0mu/}
\newcommand{\qsl}{q\mkern-9.0mu/}
\newcommand{\Asl}{A\mkern-11.0mu/}
\newcommand{\dsl}{\partial\mkern-10.0mu/}
\newcommand{\epsl}{\varepsilon\mkern-8.5mu/}
\newcommand{\vsl}{v\mkern-9.0mu/}
\newcommand{\Bsl}{B\mkern-11.0mu/}

\newcommand{\pperp}{\bm{p}_\perp}
\newcommand{\kperp}{\bm{k}_\perp}
\newcommand{\qperp}{\bm{q}_\perp}
\newcommand{\as}{\alpha_s}
\newcommand{\CF}{C_F}
\newcommand{\CA}{C_A}
\newcommand{\TF}{T_F}
\newcommand{\muMS}{\overline{\mathrm{MS}}}
\newcommand{\LQCD}{\Lambda_{\mathrm{QCD}}}
\newcommand{\Gcusp}{\Gamma_{\mathrm{cusp}}}
\newcommand{\SCETI}{\text{SCET}_{\mathrm{I}}}
\newcommand{\SCETII}{\text{SCET}_{\mathrm{II}}}
\newcommand{\en}[1]{\textsf{#1}}   % 英文术语
\newcommand{\Tr}{\mathrm{Tr}}
\newcommand{\Pexp}{\mathbf{P}\exp}

\allowdisplaybreaks[3]
\setlist{itemsep=2pt,parsep=2pt,topsep=3pt}

% =====================================================================
\title{\bfseries 软共线有效理论（SCET）课前预习讲义\\[8pt]
       \large 面向 Matthias Neubert 教授课程
       《SCET / Soft-Collinear Effective Theory》}
\author{课前自学材料（简体中文，英文术语另列）}
\date{\today}

\begin{document}
\maketitle

\begin{abstract}
\noindent
本讲义是为选修 Matthias Neubert 教授《软共线有效理论》（\en{Soft-Collinear Effective Theory}，
简称 SCET）课程的同学准备的\textbf{课前预习}材料。预设读者已经学过量子场论（正则量子化、
Feynman 规则、单圈重整化）与基本粒子物理，但\textbf{没有系统学习过 SCET}，也可能对
有效场论（\en{effective field theory}, EFT）只有零散印象。

讲义的写法遵循三条原则：（一）从相对论运动学与 EFT 的一般思想出发，不假定读者已经熟悉
光锥坐标或标度计数；（二）\textbf{完整给出推导}，尤其是光锥坐标、模式（\en{mode}）标度赋值、
以及至少一个匹配／标度计数的算例，每一步都可以自己用纸笔复核；（三）强调\textbf{物理图像}，
即“为什么大对数会出现”“软与共线到底差在哪里”，而不是堆砌高级技巧。

所有英文专业术语在正文中首次出现时用括号标注，并在第 \ref{sec:glossary} 节汇总为英汉术语表。
\end{abstract}

\vspace{0.4em}
\tableofcontents
\vspace{1.2em}

% =====================================================================
\section*{导读：如何使用这份讲义}
\addcontentsline{toc}{section}{导读：如何使用这份讲义}
\label{sec:guide}

\subsection*{这份讲义想解决的问题}

SCET 是一套用来处理\textbf{同时存在多个能量标度}、并且低能自由度\textbf{沿着光锥方向高度
不对称}的强相互作用过程的有效场论。它的技术门槛主要来自三件事：

\begin{enumerate}[label=(\arabic*)]
  \item 记号：光锥坐标、$n$ 与 $\bar n$ 参考向量、$p^+$／$p^-$／$p_\perp$、标签动量
        （\en{label momentum}）与残余动量（\en{residual momentum}）。
  \item 标度计数（\en{power counting}）：为什么共线夸克场是 $\mathcal O(\lambda)$、
        共线胶子场的三个分量各自是 $\mathcal O(\lambda^2,1,\lambda)$。
  \item 概念：$\SCETI$ 与 $\SCETII$ 的区别、软与共线的重叠（\en{zero-bin}）、
        快度发散（\en{rapidity divergence}）。
\end{enumerate}

如果第一节课就要同时消化这三样，多数人会卡在记号上而错过物理。因此本讲义把它们\textbf{拆开}：
第 \ref{sec:prelim}--\ref{sec:limits} 节只用普通的四动量与 QCD Feynman 规则，把“为什么需要
SCET”讲清楚；第 \ref{sec:lightcone} 节才引入光锥坐标；第 \ref{sec:powercounting} 节做标度
计数；第 \ref{sec:structure} 节才写下有效拉氏量。

\subsection*{建议的阅读路径}

\begin{itemize}
  \item \textbf{时间紧（约 2 小时）}：读第 \ref{sec:multiscale} 节（多标度问题）、
        第 \ref{sec:lightcone} 节（光锥坐标，务必自己推一遍）、
        第 \ref{sec:powercounting} 节的表 \ref{tab:modes}（模式标度表）、
        第 \ref{sec:glossary} 节术语表、第 \ref{sec:focus} 节听课重点。
  \item \textbf{时间充裕（约 6--8 小时）}：按顺序全读，并动手完成正文中标注的
        \textbf{练习}。特别是第 \ref{sec:eft} 节的 Fermi 理论匹配、
        第 \ref{sec:limits} 节的软胶子发射（\en{eikonal}）计算、
        第 \ref{sec:powercounting} 节的“按区域展开”（\en{expansion by regions}）算例。
  \item \textbf{只想补记号}：第 \ref{sec:lightcone} 节 $+$ 附录 \ref{app:identities}
        （Dirac 代数与光锥恒等式速查）。
\end{itemize}

\subsection*{符号约定}

全篇采用自然单位 $\hbar=c=1$，度规
\begin{equation}
  g^{\mu\nu}=\mathrm{diag}(+1,-1,-1,-1),
  \qquad
  a\cdot b = a^0b^0-\bm a\cdot\bm b .
\end{equation}
QCD 规范群为 $SU(N_c)$，$N_c=3$，$\CF=(N_c^2-1)/(2N_c)=4/3$，$\CA=N_c=3$，
$\TF=1/2$，$n_f$ 为活跃味数。维数正规化取 $d=4-2\epsilon$，减除方案为 $\muMS$。
小参数一律记作 $\lambda\ll1$。字母 $n,\bar n$ 保留给光锥参考向量，
$Q$ 表示过程的硬标度（\en{hard scale}）。

\begin{remark}
本讲义中若某个系数只需要“量级”而不需要精确值，会明确写成 $\mathcal O(\cdot)$ 或
“$\sim$”；凡是给出确定数值系数的公式，都可以（并且建议）自己验算。
少数用于定向的文献编号见第 \ref{sec:further} 节，建议听课时以老师给出的版本为准。
\end{remark}

% =====================================================================
\section{预备知识：运动学、Feynman 规则、跑动耦合}
\label{sec:prelim}

本节把后面反复使用的基础工具集中复习一遍。已经很熟的同学可以只看
第 \ref{subsec:running} 小节末尾的“大对数”讨论。

\subsection{相对论运动学：不变量、快度、相空间}
\label{subsec:kinematics}

\paragraph{Mandelstam 变量。}
对 $2\to2$ 过程 $p_1+p_2\to p_3+p_4$，定义
\begin{equation}
  s=(p_1+p_2)^2,\qquad t=(p_1-p_3)^2,\qquad u=(p_1-p_4)^2,
  \qquad s+t+u=\sum_{i=1}^{4}m_i^2 .
\end{equation}
在质心系（\en{center-of-mass frame}）中 $\sqrt s$ 就是总能量。对无质量粒子，
$s=2p_1\cdot p_2$。整个 SCET 的出发点是：\textbf{一个过程往往同时含有大不变量
（如 $s\sim Q^2$）和小不变量（如某个喷注的不变质量平方 $m_J^2\ll Q^2$）}。

\paragraph{快度与横动量。}
取 $z$ 轴为束流／喷注轴，定义横动量
$\pperp=(p^1,p^2)$、$p_T=|\pperp|$，以及快度（\en{rapidity}）
\begin{equation}
  y=\frac12\ln\frac{p^0+p^3}{p^0-p^3} .
  \label{eq:rapidity}
\end{equation}
沿 $z$ 轴的 Lorentz 推动（\en{boost}）参数 $\eta$ 作用为
\begin{equation}
  y \;\longrightarrow\; y+\eta ,
  \qquad
  \pperp \;\longrightarrow\; \pperp .
\end{equation}
即：\textbf{快度在纵向推动下只是平移，横动量不变}。这条性质是后面
“不同模式按快度分类”的根源；对无质量粒子，快度退化为伪快度
$\eta=-\ln\tan(\theta/2)$，其中 $\theta$ 是与 $z$ 轴的夹角。

小角度极限下（$\theta\ll1$）
\begin{equation}
  \eta \simeq \ln\frac{2}{\theta}\;\longrightarrow\;+\infty ,
\end{equation}
所以\textbf{共线极限就是快度发散的极限}。请记住这句话，它在第 \ref{sec:limits}
与第 \ref{sec:confusions} 节还会回来。

\paragraph{无质量两体相空间。}
单粒子相空间元
\begin{equation}
  \int\!\frac{d^3k}{(2\pi)^3 2E_k}
  =\int\!\frac{d^4k}{(2\pi)^4}\,2\pi\,\delta_+(k^2)
  =\frac{1}{(2\pi)^3}\int_0^\infty \frac{E_k\,dE_k}{2}
    \int_{-1}^{1}d\cos\theta\int_0^{2\pi}\!d\varphi ,
  \label{eq:1ps}
\end{equation}
其中 $\delta_+$ 表示只取正能解。式 \eqref{eq:1ps} 在第 \ref{subsec:eikonal}
小节推导软胶子发射时会直接用到。

\subsection{QCD Feynman 规则回顾（分层次）}
\label{subsec:feynman}

QCD 拉氏量（含规范固定项，Feynman 规范 $\xi=1$）
\begin{equation}
  \mathcal L_{\rm QCD}
  =\bar\psi\,(i\Dsl-m)\,\psi
   -\frac14 G^a_{\mu\nu}G^{a\,\mu\nu}
   -\frac{1}{2\xi}\big(\partial^\mu A^a_\mu\big)^2
   +\mathcal L_{\rm ghost},
\end{equation}
\begin{equation}
  D_\mu=\partial_\mu-ig_s A^a_\mu t^a ,
  \qquad
  G^a_{\mu\nu}=\partial_\mu A^a_\nu-\partial_\nu A^a_\mu
                +g_s f^{abc}A^b_\mu A^c_\nu .
\end{equation}
本讲义使用 $D_\mu=\partial_\mu-ig_sA_\mu$ 的约定（$A_\mu\equiv A_\mu^a t^a$）。

按“听课时最常用”的顺序，把规则分成三层：

\begin{description}[leftmargin=1.2cm,style=nextline]
\item[第一层（必须秒答）]
  夸克传播子 $\dfrac{i(\psl+m)}{p^2-m^2+i0}$；
  胶子传播子 $\dfrac{-i g^{\mu\nu}\delta^{ab}}{p^2+i0}$（Feynman 规范）；
  夸克--胶子顶点 $ig_s\gamma^\mu t^a$。
\item[第二层（会用即可）]
  三胶子、四胶子顶点；ghost 顶点；颜色因子恒等式
  \begin{equation}
    t^at^a=\CF\mathbf 1,\qquad
    \Tr[t^at^b]=\TF\delta^{ab},\qquad
    f^{acd}f^{bcd}=\CA\delta^{ab} .
  \end{equation}
\item[第三层（SCET 里会被“重写”）]
  外线极化求和、$\gamma$ 矩阵迹技术、维数正规化中的
  $\gamma^\mu\gamma_\mu=d$。SCET 会把这些量在光锥基底中重新组织，
  见附录 \ref{app:identities}。
\end{description}

\begin{keypoint}
SCET 并不引入新的 Feynman 规则“魔法”：它的传播子与顶点都是 QCD 规则在
特定动量区域中的\textbf{展开结果}。所以只要熟练掌握第一层规则，加上
一次仔细的展开，就能自己推出 SCET 的 Feynman 规则（第 \ref{subsec:scetrules} 小节）。
\end{keypoint}

\subsection{重整化与跑动：大对数从哪里来}
\label{subsec:running}

\paragraph{跑动耦合。}
维数正规化 $\muMS$ 方案中，强耦合 $\as(\mu)$ 满足
\begin{equation}
  \frac{d\as(\mu)}{d\ln\mu}
  =-\frac{\beta_0}{2\pi}\as^2+\mathcal O(\as^3),
  \qquad
  \beta_0=11-\frac23 n_f .
  \label{eq:beta}
\end{equation}
积分一次（单圈）
\begin{equation}
  \as(\mu)=\frac{\as(\mu_0)}
  {1+\dfrac{\beta_0\as(\mu_0)}{2\pi}\ln\dfrac{\mu}{\mu_0}} .
  \label{eq:alphasrun}
\end{equation}
\eqref{eq:alphasrun} 的分母里出现了 $\as\ln(\mu/\mu_0)$：这就是最简单的
“\textbf{重整化群把一串对数重求和}”的例子——展开分母得到
$\as\sum_n\big(-\tfrac{\beta_0\as}{2\pi}L\big)^n$，
即把所有 $\as^{n+1}L^n$ 一次性收拢。

\paragraph{反常量纲的作用。}
更一般地，若某个（重整化后的）系数 $C(\mu)$ 满足
\begin{equation}
  \frac{dC(\mu)}{d\ln\mu}=\gamma(\as)\,C(\mu),
\end{equation}
则
\begin{equation}
  C(\mu)=C(\mu_0)\,\exp\!\left[\int_{\ln\mu_0}^{\ln\mu}\!
  \gamma\big(\as(\mu')\big)\,d\ln\mu'\right].
  \label{eq:rgesol}
\end{equation}
把 $C(\mu_0)$ 在 $\mu_0\sim$（该系数的“自然标度”）处用固定阶
微扰论计算，再用 \eqref{eq:rgesol} 演化到我们真正需要的标度 $\mu$，
就得到\textbf{重求和}（\en{resummation}）的结果。这正是 SCET 的最终用途。

\paragraph{两种对数：单对数与双对数。}
上面 $\as^n L^n$ 的结构叫\textbf{单对数}（\en{single logarithm}），
来自紫外／共线的一重发散。而在含有软与共线两类发散的过程中，
每阶 $\as$ 最多带来\textbf{两个}对数：
\begin{equation}
  \text{幅度或截面} \;\sim\; \sum_n \as^n
  \big(c_{n,2n}L^{2n}+c_{n,2n-1}L^{2n-1}+\dots\big),
  \qquad L=\ln\frac{Q}{\mu_{\rm soft}} .
  \label{eq:doublelogs}
\end{equation}
这类 $\as^nL^{2n}$ 称为\textbf{Sudakov 双对数}（\en{Sudakov double logarithm}）。
当 $\as L^2\sim1$ 时固定阶微扰论彻底失效，即使 $\as\ll1$。

\begin{keypoint}
一句话总结整门课的动机：\textbf{当过程中出现层级化的标度
$Q\gg\mu_1\gg\mu_2\gg\dots$ 时，固定阶微扰论里会出现 $\as^nL^{2n}$；
SCET 的作用是把每个标度装进一个独立的因子（硬函数、喷注函数、软函数），
让每个因子在自己的自然标度上无大对数，再用重整化群把它们连接起来。}
\end{keypoint}

\begin{exercise}
用 \eqref{eq:beta} 验证 \eqref{eq:alphasrun}；再证明用 $\mu^2$ 作变量时
$d\as/d\ln\mu^2=-\beta_0\as^2/(4\pi)$。注意区分 $\ln\mu$ 与 $\ln\mu^2$
的约定，SCET 文献中两者都常见，差一个因子 2 是最常见的低级错误。
\end{exercise}

% =====================================================================
\section{散射运动学与多标度问题：为什么需要有效理论}
\label{sec:multiscale}

\subsection{一个具体的“坏”例子：$e^+e^-\to$ 喷注的 thrust 分布}
\label{subsec:thrustintro}

考虑 $e^+e^-$ 在质心能量 $Q$ 上湮灭产生强子。定义 thrust（\en{thrust}）
\begin{equation}
  T=\max_{\bm n}\frac{\sum_i|\bm p_i\cdot\bm n|}{\sum_i|\bm p_i|},
  \qquad \tau\equiv 1-T .
  \label{eq:thrustdef}
\end{equation}
$\tau\to0$ 对应两个极窄的背对背喷注（\en{back-to-back jets}），
$\tau$ 越大事件越“胖”。在双喷注极限下有一个很有用的等价表述
\begin{equation}
  \tau\;\simeq\;\frac{m_L^2+m_R^2}{Q^2},
  \label{eq:thrusthemisphere}
\end{equation}
其中 $m_{L,R}$ 是把事件按 thrust 轴分成两个半球后各半球的总不变质量。

用固定阶 QCD 计算 $\mathcal O(\as)$ 的 thrust 分布，在 $\tau\to0$ 的
奇异部分（\en{singular part}）为
\begin{equation}
  \frac{1}{\sigma_0}\frac{d\sigma}{d\tau}
  \;\xrightarrow[\ \tau\to0\ ]{}\;
  \frac{\as \CF}{2\pi}\,\frac{1}{\tau}
  \Big(-4\ln\tau-3\Big)+\dots
  \label{eq:thrustfixed}
\end{equation}
（$\sigma_0$ 为 Born 截面，$\dots$ 表示 $\tau\to0$ 时不奇异的项与 $\delta(\tau)$ 项）。
把它累积起来：
\begin{equation}
  \frac{\sigma(\tau)}{\sigma_0}
  =\int_0^\tau\! d\tau'\,\frac{1}{\sigma_0}\frac{d\sigma}{d\tau'}
  =1+\frac{\as\CF}{2\pi}\Big(-2\ln^2\tau-3\ln\tau\Big)+\mathcal O(\as^2).
  \label{eq:thrustcumul}
\end{equation}

\eqref{eq:thrustcumul} 就是 \eqref{eq:doublelogs} 的具体实现：
系数是 $\as\ln^2\tau$。取 $Q=m_Z\simeq91\,$GeV、$\as\simeq0.12$，
当 $\tau\sim10^{-3}$ 时 $\as\ln^2\tau\sim0.12\times(6.9)^2/(2\pi)\times\tfrac43\times2\approx 2$，
微扰级数已经不收敛。而实验上恰恰是小 $\tau$ 区域统计量最大、最适合精确测
$\as$。这就是必须重求和的现实压力。

\subsection{标度层级：一个过程里到底有几个标度}
\label{subsec:hierarchy}

回到 \eqref{eq:thrusthemisphere}。在双喷注构型下：

\begin{itemize}
  \item 两个喷注的\textbf{总能量}是 $\mathcal O(Q)$——这是\textbf{硬标度}
        （\en{hard scale}）。
  \item 每个喷注的不变质量满足 $m_J^2\sim \tau Q^2$，即
        $m_J\sim Q\sqrt\tau$——这是\textbf{喷注标度}（\en{jet scale}）。
  \item 喷注之间大角度交换的软辐射，其能量 $\sim Q\tau$——这是
        \textbf{软标度}（\en{soft scale}）。因为半球质量
        $m^2 \supset 2E_{\rm jet}\,k^0(1-\cos\theta)\sim Q\,k^0$，
        要贡献 $m^2\sim\tau Q^2$ 就需要 $k^0\sim\tau Q$。
\end{itemize}

于是
\begin{equation}
  \boxed{\;Q\;\gg\;Q\sqrt\tau\;\gg\;Q\tau\;}
  \qquad(\tau\ll1),
\end{equation}
三个标度呈等比层级。表 \ref{tab:processes} 列出若干常见过程的标度层级。

\begin{table}[htbp]
\centering
\caption{几个典型过程中的标度层级。$\lambda$ 为 SCET 展开小参数；
最后一列的类型见第 \ref{subsec:scetI_II} 小节。}
\label{tab:processes}
\begin{tabular}{@{}llll@{}}
\toprule
过程 / 可观测量 & 硬标度 & 中间（共线）标度 & 软标度 \\
\midrule
$e^+e^-\to$ 2 jets, thrust $\tau$ & $Q$ & $Q\sqrt\tau$ & $Q\tau$ \\
DIS，$x\to1$                      & $Q$ & $Q\sqrt{1-x}$ & $Q(1-x)$ \\
$B\to X_s\gamma$ 端点区           & $m_b$ & $\sqrt{m_b\LQCD}$ & $\LQCD$ \\
Drell--Yan，小 $q_T$              & $Q$ & $\sqrt{Q q_T}$ 或 $q_T$ & $q_T$ \\
$B\to\pi\ell\nu$（$\SCETII$）      & $m_b$ & $\LQCD$ & $\LQCD$ \\
\bottomrule
\end{tabular}
\end{table}

\subsection{为什么“直接算”不行}

有人会问：既然固定阶 QCD 原则上给出正确答案，为什么不多算几阶？三个理由：

\begin{enumerate}[label=(\roman*)]
  \item \textbf{收敛性}：如上所述，$\as L^2\sim1$ 时每一阶都同等重要，
        加阶数不解决问题，必须\textbf{对全部阶重求和}。
  \item \textbf{结构性}：重求和需要知道对数系数的\textbf{全阶结构}。
        EFT 提供的因子化定理（\en{factorization theorem}）恰恰把这个结构
        变成“每个因子满足一个重整化群方程”，从而可以指数化。
  \item \textbf{非微扰输入}：当软标度降到 $\LQCD$ 附近时，软函数
        不再是微扰量。EFT 给出\textbf{清晰的算符定义}，让非微扰部分
        被隔离成一个普适的、可以从实验或格点提取的对象
        （如形状函数 \en{shape function}、光锥分布振幅 \en{light-cone distribution amplitude}）。
\end{enumerate}

\begin{keypoint}
EFT 的价值不只是“算得更准”，更重要的是\textbf{把不同标度的物理分离成独立、
可分别定义、可分别计算或测量的对象}。这一点在第 \ref{sec:eft} 节展开。
\end{keypoint}

\begin{exercise}
由 \eqref{eq:thrustfixed} 出发验证 \eqref{eq:thrustcumul}
（提示：$\int_0^\tau d\tau'\,(-4\ln\tau'-3)/\tau'$ 在 $\tau'\to0$ 端点发散，
需要与 $\delta(\tau)$ 处的虚修正相加才有限；这正是“加号分布”
\en{plus distribution} 的用途，见第 \ref{sec:confusions} 节第 \ref{item:plusdist} 条）。
\end{exercise}

% =====================================================================
\section{有效场论总览：积分掉重模、匹配、Wilson 系数}
\label{sec:eft}

\subsection{EFT 的三步法}

任何一个有效场论的构造都可以概括为三步：

\begin{enumerate}[label=\textbf{第 \arabic* 步}, leftmargin=2.1cm]
  \item \textbf{识别自由度与展开参数}。写下在所关心的能量区域里\textbf{仍然
        动力学活跃}的场，以及被“积分掉”（\en{integrated out}）的重自由度，
        确定小参数（$\lambda=$ 轻标度／重标度）。
  \item \textbf{写下最一般的有效拉氏量}。按 $\lambda$ 的幂次组织算符，
        只受对称性约束：
        \begin{equation}
          \mathcal L_{\rm eff}=\sum_{k}\sum_{i} C_i^{(k)}(\mu)\,
          \frac{O_i^{(k)}}{M^{k}} .
          \label{eq:leff}
        \end{equation}
        算符 $O_i$ 只含轻自由度；重物理全部藏进系数 $C_i$。
  \item \textbf{匹配（\en{matching}）与跑动}。在某个标度 $\mu\sim M$ 处要求
        全理论与有效理论的（在轻标度上展开后的）$S$ 矩阵元
        \textbf{逐阶相等}，解出 Wilson 系数（\en{Wilson coefficient}）$C_i(\mu\sim M)$；
        再用重整化群把 $C_i$ 演化到 $\mu\sim$ 轻标度。
\end{enumerate}

\begin{keypoint}
匹配的本质：\textbf{全理论与有效理论的紫外行为不同，但红外行为必须相同}。
两者之差只含短距离信息，因此是 $\LQCD$ 或轻标度的解析函数，可以吸收进
局域算符的系数。这句话在 SCET 中依然成立，只不过“局域”要放宽为
“沿光锥方向非局域”（见第 \ref{sec:structure} 节）。
\end{keypoint}

\subsection{范例一：Fermi 理论的完整匹配（含推导）}
\label{subsec:fermi}

这是最干净的匹配算例，请务必自己推一遍。考虑 $b\to c\bar u d$ 的树图，
由 $W$ 玻色子交换给出。全理论振幅：
\begin{equation}
  i\mathcal M_{\rm full}
  =\left(\frac{-ig}{2\sqrt2}\right)^2 V_{cb}V_{ud}^*
  \Big[\bar u_c\gamma^\mu(1-\gamma_5)u_b\Big]
  \Big[\bar u_d\gamma^\nu(1-\gamma_5)u_u\Big]
  \frac{-i\left(g_{\mu\nu}-\dfrac{q_\mu q_\nu}{M_W^2}\right)}{q^2-M_W^2} ,
  \label{eq:fullamp}
\end{equation}
其中 $q$ 是 $W$ 传递的动量，$q^2\sim m_b^2\ll M_W^2$。

\paragraph{展开传播子。}这是“积分掉重模”的字面含义：
\begin{equation}
  \frac{1}{q^2-M_W^2}
  =-\frac{1}{M_W^2}\cdot\frac{1}{1-q^2/M_W^2}
  =-\frac{1}{M_W^2}\left(1+\frac{q^2}{M_W^2}+\frac{q^4}{M_W^4}+\dots\right).
  \label{eq:propexpand}
\end{equation}
注意 $q_\mu q_\nu/M_W^2$ 项在与流收缩后给出正比于夸克质量的贡献
（用运动方程 $\qsl\to m$），也被 $1/M_W^2$ 压低。保留领头项：
\begin{equation}
  i\mathcal M_{\rm full}
  = -i\,\frac{g^2}{8M_W^2}\,V_{cb}V_{ud}^*
    \Big[\bar u_c\gamma^\mu(1-\gamma_5)u_b\Big]
    \Big[\bar u_d\gamma_\mu(1-\gamma_5)u_u\Big]
    +\mathcal O\!\left(\frac{q^2}{M_W^4}\right).
  \label{eq:fullexpanded}
\end{equation}

\paragraph{有效理论一侧。}低能理论中不再有 $W$ 场，只有四费米子局域算符
\begin{equation}
  \mathcal L_{\rm eff}
  =-\frac{4G_F}{\sqrt2}V_{cb}V_{ud}^*\,
   \Big[C_1(\mu)\,O_1+C_2(\mu)\,O_2\Big]+\text{h.c.}+\mathcal O(1/M_W^4),
  \label{eq:lefffermi}
\end{equation}
\begin{equation}
  O_1=\big[\bar c\,\gamma^\mu P_L\, b\big]\big[\bar d\,\gamma_\mu P_L\, u\big],
  \qquad
  O_2=\big[\bar c\,\gamma^\mu P_L\, u\big]\big[\bar d\,\gamma_\mu P_L\, b\big],
  \qquad P_L=\frac{1-\gamma_5}{2}.
\end{equation}
（$O_2$ 是 $O_1$ 的“颜色重排”版本，树图上不产生，单圈胶子交换后必须引入。）

\paragraph{匹配条件。}比较 \eqref{eq:fullexpanded} 与
\eqref{eq:lefffermi} 的树级矩阵元，得到
\begin{equation}
  \boxed{\;\frac{G_F}{\sqrt2}=\frac{g^2}{8M_W^2}\;},
  \qquad
  C_1(\mu\!\sim\! M_W)=1+\mathcal O(\as),
  \qquad
  C_2(\mu\!\sim\! M_W)=0+\mathcal O(\as).
  \label{eq:fermimatch}
\end{equation}

\paragraph{跑动。}到 $\mathcal O(\as)$，$O_{1,2}$ 会\textbf{混合}，
其 Wilson 系数满足矩阵形式的重整化群方程
\begin{equation}
  \frac{d}{d\ln\mu}\begin{pmatrix}C_1\\C_2\end{pmatrix}
  =\hat\gamma^{T}(\as)\begin{pmatrix}C_1\\C_2\end{pmatrix},
  \qquad
  \hat\gamma=\frac{\as}{4\pi}
  \begin{pmatrix}-6/N_c & 6\\ 6 & -6/N_c\end{pmatrix}+\mathcal O(\as^2).
  \label{eq:admatrix}
\end{equation}
把 $C_i$ 从 $\mu=M_W$ 演化到 $\mu=m_b$，就把所有
$\as^n\ln^n(M_W/m_b)$ 重求和了。

\begin{keypoint}
Fermi 理论示范了 EFT 的全部要素，但它只有\textbf{单对数}
$\as^n\ln^n(M_W/m_b)$：这是因为被积掉的重模与保留的轻模在\textbf{虚度}
（\en{virtuality}）上一刀两断，红外结构简单。SCET 的困难在于低能理论里
\textbf{同时存在几种虚度／快度都不同的轻模}，从而出现双对数。
\end{keypoint}

\subsection{范例二：HQET —— 场的重新定义与标度赋值的预演}
\label{subsec:hqet}

重夸克有效理论（\en{Heavy Quark Effective Theory}, HQET）是 SCET 的直接前身，
Neubert 本人是该领域的奠基者之一，课上极可能作为类比反复引用。

重夸克动量分解为“大部分 $+$ 残余”：
\begin{equation}
  p^\mu = m_Q\, v^\mu + k^\mu ,
  \qquad v^2=1,\qquad k^\mu\sim\LQCD\ll m_Q .
  \label{eq:hqetmom}
\end{equation}
定义把快速振荡相位剥离后的场
\begin{equation}
  h_v(x)=e^{\,i m_Q v\cdot x}\,\frac{1+\vsl}{2}\,Q(x),
  \qquad \vsl\equiv v_\mu\gamma^\mu ,
\end{equation}
则领头阶拉氏量为
\begin{equation}
  \mathcal L_{\rm HQET}
  =\bar h_v\,(i v\cdot D)\,h_v
   +\frac{1}{2m_Q}\Big[\bar h_v (iD_\perp)^2 h_v
   +\frac{c_{\rm mag}(\mu)}{2}\,\bar h_v\, g_s\sigma_{\mu\nu}G^{\mu\nu} h_v\Big]
   +\mathcal O\!\left(\frac{1}{m_Q^2}\right).
  \label{eq:lhqet}
\end{equation}

请注意其中三个将在 SCET 里原封不动复用的思想：

\begin{enumerate}[label=(\alph*)]
  \item \textbf{相位剥离}：$e^{imv\cdot x}$ 把“大动量”变成场的\textbf{标签}，
        剩下的场只携带小动量 $k$。SCET 的标签动量（\en{label momentum}）
        就是这个想法的推广。
  \item \textbf{投影算符}：$(1+\vsl)/2$ 只保留粒子解，把反粒子（大分量）积掉。
        SCET 中对应的是光锥投影 $\nsl\nbsl/4$。
  \item \textbf{按 $1/m$ 展开的算符阶梯}：领头阶极简，高阶项系数需要匹配。
        SCET 中的展开参数是 $\lambda$，结构完全类似。
\end{enumerate}

\begin{exercise}
从 QCD 的 $\bar Q(i\Dsl-m_Q)Q$ 出发，代入
$Q(x)=e^{-im_Qv\cdot x}\big[h_v(x)+H_v(x)\big]$，
其中 $H_v=e^{im_Qv\cdot x}\tfrac{1-\vsl}{2}Q$，
用 $H_v$ 的运动方程消去 $H_v$，验证 \eqref{eq:lhqet} 的前两项。
这个练习与第 \ref{subsec:derivL} 小节把 SCET 拉氏量推出来的步骤\textbf{一模一样}，
先做它会让后面轻松很多。
\end{exercise}

% =====================================================================
\section{软极限与共线极限：物理图像与奇异性来源}
\label{sec:limits}

本节回答一个关键问题：\textbf{QCD 振幅在什么情况下发散，发散来自哪个传播子？}
答案有两个（也只有两个，这是 QCD 的红外结构定理内容）：
\textbf{软}（\en{soft}）与\textbf{共线}（\en{collinear}）。

\subsection{软极限与 eikonal 近似（完整推导）}
\label{subsec:eikonal}

设有一个硬过程，其中有一条出射的无质量夸克线，动量 $p$（$p^2=0$）。
现在从这条腿上发射一个软胶子，动量 $k$、颜色 $a$、极化 $\varepsilon$，
且 $k^\mu\to0$。修改后的振幅为
\begin{equation}
  \mathcal M^{\mu}
  =\bar u(p)\,\big(ig_s\gamma^\mu t^a\big)\,
   \frac{i\big(\psl+\ksl\big)}{(p+k)^2+i0}\,\Gamma\,\big|\dots\big\rangle ,
  \label{eq:softamp1}
\end{equation}
其中 $\Gamma$ 代表硬过程其余部分。分母：
\begin{equation}
  (p+k)^2=p^2+2p\cdot k+k^2=2p\cdot k+k^2\;\xrightarrow[\ k\to0\ ]{}\;2p\cdot k .
\end{equation}
分子用 Dirac 代数处理。利用
$\gamma^\mu\psl=2p^\mu-\psl\gamma^\mu$ 与外线运动方程 $\bar u(p)\psl=0$：
\begin{align}
  \bar u(p)\,\gamma^\mu(\psl+\ksl)
  &=\bar u(p)\big(2p^\mu-\psl\gamma^\mu\big)+\bar u(p)\gamma^\mu\ksl
   \notag\\
  &=2p^\mu\,\bar u(p)+\underbrace{\bar u(p)\gamma^\mu\ksl}_{\mathcal O(k)} .
  \label{eq:eiknum}
\end{align}
在 $k\to0$ 下第二项相对第一项被压低一个 $k$ 的幂次，可以丢掉。于是
\begin{equation}
  \boxed{\;
  \mathcal M^\mu \;\simeq\;
  g_s\,t^a\,\frac{p^\mu}{p\cdot k}\;\mathcal M_0\;}
  \label{eq:eikonal}
\end{equation}
（$\mathcal M_0$ 为不发射软胶子的振幅）。这就是\textbf{eikonal 近似}
（\en{eikonal approximation}，也叫 soft approximation）。

\begin{remark}
\eqref{eq:eikonal} 的物理含义极其重要：软胶子\textbf{完全看不见}发射它的粒子的
自旋，也看不见硬过程的内部结构，只“看到”粒子的\textbf{方向与颜色荷}。
这就是为什么软辐射可以用 Wilson 线（\en{Wilson line}）——即经典色荷沿直线运动
产生的规范场——来完整描述。$p^\mu/(p\cdot k)$ 正是 Wilson 线的 Feynman 规则。
\end{remark}

\subsection{软发射给出双对数：显式计算}
\label{subsec:doublelog}

现在把 \eqref{eq:eikonal} 用于 $e^+e^-\to q\bar q$ 加一个软胶子。
取质心系
\begin{equation}
  p_1=E(1,0,0,1),\qquad p_2=E(1,0,0,-1),\qquad Q=2E ,
\end{equation}
\begin{equation}
  k=E_k\big(1,\sin\theta\cos\varphi,\sin\theta\sin\varphi,\cos\theta\big).
\end{equation}
两条腿都可以发射，总 eikonal 电流为
$J^\mu=g_st^a\big(\frac{p_1^\mu}{p_1\cdot k}-\frac{p_2^\mu}{p_2\cdot k}\big)$
（相对符号来自出射夸克与出射反夸克的色荷相反）。对极化与颜色求和：
\begin{align}
  \sum_{\rm pol,col}|J|^2
  &= g_s^2\,\CF\left[-\left(\frac{p_1}{p_1\cdot k}-\frac{p_2}{p_2\cdot k}\right)^2\right]
   = g_s^2\,\CF\,\frac{2\,p_1\cdot p_2}{(p_1\cdot k)(p_2\cdot k)} ,
  \label{eq:eiksq}
\end{align}
这里用了 $p_{1,2}^2=0$。代入运动学：
\begin{equation}
  p_1\cdot k=E E_k(1-\cos\theta),\qquad
  p_2\cdot k=E E_k(1+\cos\theta),\qquad
  2p_1\cdot p_2=4E^2 ,
\end{equation}
\begin{equation}
  \Rightarrow\quad
  \frac{2p_1\cdot p_2}{(p_1\cdot k)(p_2\cdot k)}
  =\frac{4E^2}{E^2E_k^2(1-\cos^2\theta)}
  =\frac{4}{E_k^2\sin^2\theta} .
\end{equation}
配上相空间 \eqref{eq:1ps}：
\begin{align}
  \frac{d\sigma}{\sigma_0}
  &= g_s^2\CF\,\frac{1}{(2\pi)^3}\,\frac{E_k\,dE_k}{2}\,
     d\cos\theta\, d\varphi\;\frac{4}{E_k^2\sin^2\theta}
   \notag\\
  &= \frac{4\pi\as\,\CF\cdot 4}{2(2\pi)^3}\,
     \frac{dE_k}{E_k}\,\frac{d\cos\theta}{\sin^2\theta}\,d\varphi
   \notag\\
  &= \frac{2\as\CF}{\pi}\,\frac{dE_k}{E_k}\,
     \frac{d\cos\theta}{\sin^2\theta} ,
  \label{eq:softdist}
\end{align}
最后一步用了 $g_s^2=4\pi\as$ 与 $\int_0^{2\pi}d\varphi=2\pi$。

现在读出奇异性。在 $\theta\to0$ 附近
$|d\cos\theta|=\sin\theta\,d\theta$、$\sin^2\theta\simeq\theta^2$，所以
\begin{equation}
  \boxed{\;
  \frac{d\sigma}{\sigma_0}\;\simeq\;
  \frac{2\as\CF}{\pi}\,\frac{dE_k}{E_k}\,\frac{d\theta}{\theta}\;}
  \qquad(\theta\to0\ \text{或}\ \theta\to\pi).
  \label{eq:doublelogdiff}
\end{equation}

两个独立的对数发散清清楚楚：

\begin{itemize}
  \item $dE_k/E_k$：\textbf{软发散}（\en{soft divergence}），来自
        $E_k\to0$，即 \eqref{eq:eikonal} 分母中的 $p\cdot k\to0$
        因为 $k\to0$。
  \item $d\theta/\theta$：\textbf{共线发散}（\en{collinear divergence}），来自
        $\theta\to0$，即 $p\cdot k\to0$ 因为 $k$ 与 $p$ 平行
        （即便 $E_k$ 不小！）。
\end{itemize}

若把积分限设为 $\mu_{\rm IR}<E_k<Q$、$\theta_0<\theta<1$，则
\begin{equation}
  \frac{\Delta\sigma}{\sigma_0}\sim
  \frac{2\as\CF}{\pi}\ln\frac{Q}{\mu_{\rm IR}}\ln\frac{1}{\theta_0}
  \;\sim\;\as L^2 .
\end{equation}
这就是 \eqref{eq:doublelogs} 与 \eqref{eq:thrustcumul} 中双对数的来源：
\textbf{一个来自能量，一个来自角度（等价地：快度）}。

\begin{keypoint}
\textbf{双对数 $=$ 软（能量）对数 $\times$ 共线（角度／快度）对数。}
SCET 的全部结构，本质上都是为了把这两类对数分别装进
“软函数”与“喷注函数”，从而各自只含单一标度。
\end{keypoint}

\subsection{共线极限与分裂函数}
\label{subsec:collinear}

共线发散也可以在\textbf{不取软极限}的情况下单独看到。设 $p$ 分裂为
两个近共线的动量 $p_1=z p+\dots$、$p_2=(1-z)p+\dots$，
则中间传播子的虚度为
\begin{equation}
  (p_1+p_2)^2 \simeq \frac{\kperp^2}{z(1-z)} ,
  \label{eq:splitvirt}
\end{equation}
其中 $\kperp$ 是相对横动量。$\kperp\to0$ 时传播子发散，且 $z$ 保持有限。
积掉这个区域给出著名的 Altarelli--Parisi 分裂函数（\en{splitting function}）：
\begin{equation}
  d\sigma_{n+1}\simeq d\sigma_n\,\frac{\as}{2\pi}\,
  \frac{d\kperp^2}{\kperp^2}\,dz\;P_{ji}(z),
\end{equation}
\begin{equation}
  P_{qq}(z)=\CF\,\frac{1+z^2}{1-z},\qquad
  P_{gq}(z)=\CF\,\frac{1+(1-z)^2}{z},\qquad
  P_{gg}(z)=2\CA\!\left[\frac{z}{1-z}+\frac{1-z}{z}+z(1-z)\right].
\end{equation}
注意 $P_{qq}(z)$ 在 $z\to1$ 时还有一个 $1/(1-z)$ 的极点——那正是
\textbf{软}极限，它同时是共线的：这就是软与共线区域的\textbf{重叠}
（\en{overlap}），在 SCET 里必须小心处理（\en{zero-bin subtraction}，
见第 \ref{sec:confusions} 节）。

\subsection{奇异性的来源：一句话总结}

回顾一下 \eqref{eq:softamp1}--\eqref{eq:splitvirt}，所有红外奇异性都来自
\textbf{某个内线传播子的分母趋于零}：
\begin{equation}
  \frac{1}{(p+k)^2}=\frac{1}{2p\cdot k}
  =\frac{1}{2E_pE_k(1-\cos\theta_{pk})} \;\to\;\infty
  \quad\Longleftrightarrow\quad
  \underbrace{E_k\to0}_{\text{软}}
  \ \ \text{或}\ \
  \underbrace{\theta_{pk}\to0}_{\text{共线}} .
  \label{eq:irorigin}
\end{equation}

\begin{remark}[KLN 与 IR 安全]
对足够“包含性”（\en{inclusive}）的可观测量，实发射与虚修正的红外发散
互相抵消（Kinoshita--Lee--Nauenberg 定理），末态截面有限。但\textbf{留下的
不是零，而是大对数}——因为可观测量本身（如 $\tau$）给出了一个截断。
换句话说：\textbf{测量行为把红外发散转化为对数增强}，
而这些对数的比值标度正是表 \ref{tab:processes} 中列出的那些。
\end{remark}

\begin{exercise}
用 \eqref{eq:eikonal} 证明式 \eqref{eq:eiksq} 中的等式
$-\big(\frac{p_1}{p_1\cdot k}-\frac{p_2}{p_2\cdot k}\big)^2
=\frac{2p_1\cdot p_2}{(p_1\cdot k)(p_2\cdot k)}$，
并说明为什么必须用 $\sum_{\rm pol}\varepsilon^\mu\varepsilon^{*\nu}\to-g^{\mu\nu}$
才是合法的（提示：eikonal 电流守恒 $k\cdot J=0$）。
\end{exercise}

\end{document}
